%% Section 3: The Integrable Lattice Model Perspective
%% To be \input{} into the main document.

\section{The Integrable Lattice Model Perspective}
\label{sec:lattice}

Section~\ref{sec:classical} closed with an observation: the transfer
operators $T_{\mathrm{free}}$ and $\bar{T}_{\mathrm{free}}$ are row
transfer matrices of a lattice model, and the matrix element
$\langle\varnothing|T_{\mathrm{free}}(u_1)\cdots
T_{\mathrm{free}}(u_n)|\lambda\rangle$ is a partition function---a
weighted sum over all configurations with prescribed boundary
conditions.  In this section we make the lattice model explicit.
The punchline is that the puzzle pieces of Knutson-Tao-Woodward are
\emph{Boltzmann weights} of an exactly solvable vertex model, and the
exact solvability follows from a single algebraic identity: the
Yang-Baxter equation.

The reader familiar with symmetric functions but not with statistical
mechanics should think of a vertex model as follows.  One draws a
planar graph (the lattice), assigns labels---called
\emph{spins}---to the edges, and assigns a \emph{weight} to each
vertex as a function of its surrounding spins.  A \emph{configuration}
is a global assignment of spins to all edges.  The \emph{partition
function} is the sum of the products of all vertex weights, taken over
all configurations with fixed boundary spins.  The partition function
is the quantity one wants to compute; the Yang-Baxter equation is what
makes the computation tractable.

\subsection{Zinn-Justin's observation}
\label{sec:lattice-zj}

The starting point is the paper of Zinn-Justin
\cite{zinn-justin-2008}, which recast the KTW puzzle model in the
language of integrable lattice models.  We present the key ideas in
three steps: reformulation, integrability, and consequences.

\subsubsection*{Step 1: Puzzles as a vertex model}

A KTW puzzle (Definition~\ref{def:ktw-puzzle}) tiles an equilateral
triangle with unit triangles and rhombi, each edge carrying a label
$0$ or~$1$.  The internal geometry consists of two families of
triangles---upward-pointing ($\triangle$) and downward-pointing
($\nabla$)---arranged in rows.  Each row processes a horizontal slice
of the puzzle triangle, reading edge labels from left to right.

To pass from a tiling model to a vertex model, one performs a standard
duality: each \emph{rhombus} (two triangles sharing an internal edge)
becomes a \emph{vertex}, and the labels on its four external edges
become the spins.  The three orientations of rhombus correspond to
three orientations of vertex.  But a more economical reformulation is
available: decompose the row processing into elementary steps, one
$\triangle$ or $\nabla$ at a time.  Each elementary step reads two
input edge labels and produces one output label (for $\triangle$,
which is deterministic) or reads one input label and produces two
outputs with possible branching (for $\nabla$, which is where all the
multiplicity lives).  Composing these steps across a row yields the
row transfer matrix $T(u)$.

The upshot is that a KTW puzzle of side $N$ becomes a vertex model on
an $N$-row lattice, where each row has $2i-1$ vertices (at row $i$
from the top), the boundary spins are determined by the partitions
$\lambda$, $\mu$, $\nu$, and each configuration of internal spins is
a valid puzzle tiling.  The partition function---the sum of the
products of all vertex weights---equals the LR coefficient
$c^{\nu}_{\lambda,\mu}$.

\begin{remark}
Since all nonzero vertex weights in the KTW model are equal to~$1$
(the model is \emph{unweighted}), the partition function is simply a
\emph{count} of valid configurations.  This is why the LR coefficient
is a nonneg\-ative integer: it counts puzzle tilings, hence
configurations of a vertex model with unit weights.
\end{remark}

\subsubsection*{Step 2: The Yang-Baxter equation}

The Yang-Baxter equation (YBE) is the fundamental integrability
condition.  In its graphical form, it asserts an equivalence between
two ways of routing three lines past one another:
\begin{equation}
\label{eq:ybe}
  R_{12}(u-v)\, R_{13}(u)\, R_{23}(v)
  \;=\;
  R_{23}(v)\, R_{13}(u)\, R_{12}(u-v),
\end{equation}
where $R_{ij}(u)$ is a matrix acting on the tensor product of the
$i$-th and $j$-th copies of the local spin space $\mathbb{C}^2$, and
$u, v$ are \emph{spectral parameters}---continuous deformation
variables that will play a central role in
Section~\ref{sec:hierarchy}.  Graphically, the left-hand side
represents three lines crossing in the order $(1,2), (1,3), (2,3)$
from right to left, and the right-hand side represents the opposite
order $(2,3), (1,3), (1,2)$.  The YBE states that these two composite
crossings give the same overall map.

\begin{remark}
The YBE is a system of cubic polynomial equations in the matrix
entries of $R$.  For $2 \times 2$ spin spaces, $R$ is a $4 \times 4$
matrix with $16$ entries, and the YBE imposes $64$ equations (one for
each matrix element of the $8 \times 8$ equation~\eqref{eq:ybe}).
That the six puzzle weights---most of them equal to~$1$---satisfy this
overdetermined system is a genuine surprise, not something one would
guess from the puzzle picture.  It is this ``miracle'' that
Zinn-Justin's observation brought to light.
\end{remark}

\subsubsection*{Step 3: Transfer matrix commutativity}

The YBE has a celebrated consequence: the \emph{commutativity of
transfer matrices}.  Define the row transfer matrix
\begin{equation}
\label{eq:transfer-matrix}
  T(u) \;=\; \mathrm{tr}_a\bigl(
    R_{a,1}(u)\, R_{a,2}(u) \cdots R_{a,N}(u)
  \bigr),
\end{equation}
where the trace is over an auxiliary space $a$ and the matrices
$R_{a,i}(u)$ act on the tensor product of $a$ with the $i$-th
``physical'' spin.  The transfer matrix processes one row of the
lattice: it reads the spin configuration below and produces the
configuration above, summing over the auxiliary spin.

The YBE implies that transfer matrices at different values of the
spectral parameter commute:
\begin{equation}
\label{eq:transfer-commute}
  [T(u),\, T(v)] = 0 \qquad \text{for all } u, v.
\end{equation}
The proof is a direct consequence of~\eqref{eq:ybe}: one ``pushes''
the auxiliary line of $T(u)$ past the auxiliary line of $T(v)$ using
$N$ applications of the YBE, one at each physical site.  The
R-matrices appear and cancel in pairs, leaving the two transfer
matrices in reversed order.

This is the connection to Section~\ref{sec:classical-hopf}.  The
operators $T_{\mathrm{free}}(u)$ and $\bar{T}_{\mathrm{free}}(v)$
defined there are precisely the transfer matrices of the
free-fermionic vertex model at spectral parameters $u$ and $v$.
Their commutativity---which we used to compute the coproduct
$\Delta(s_\lambda)$---is not an algebraic coincidence.  It is the
Yang-Baxter equation in action.

\subsubsection*{What integrability buys}

Without the YBE, computing the partition function $Z$ (which equals
the LR coefficient) requires enumerating all valid configurations of
the lattice model---a sum over exponentially many terms.  With the
YBE, the situation changes qualitatively.

First, the commuting family $\{T(u)\}_{u \in \mathbb{C}}$ can be
\emph{simultaneously diagonalized}.  The common eigenvectors are
constructed by the \emph{Bethe ansatz}, a technique from quantum
inverse scattering that expresses each eigenvector as a product of
creation operators applied to a reference state, with the creation
parameters (the ``Bethe roots'') determined by a system of algebraic
equations.

Second, the partition function can be expressed as a product of
eigenvalues, bypassing the configuration sum entirely.  One passes
from an \emph{enumerative} computation (count all tilings) to an
\emph{algebraic} one (diagonalise a commuting family of matrices).
This is the fundamental advantage of integrability: it replaces
brute-force enumeration with structure.

For the LR coefficient specifically, the Bethe ansatz machinery
yields determinantal formulas, contour integral representations, and
functional equations that would be invisible from the puzzle picture
alone.  The Yang-Baxter equation does not merely \emph{explain} why
puzzles work---it gives access to tools that go beyond what any
single combinatorial model can offer.

\subsection{The R-matrix}
\label{sec:lattice-rmatrix}

We now write down the R-matrix explicitly.  A vertex of the model
has four edges: two horizontal (``in'' from the left and ``out'' to
the right) and two vertical (``in'' from below and ``out'' above).
Each edge carries a spin in $\{0, 1\}$.  The R-matrix
$R \colon \mathbb{C}^2 \otimes \mathbb{C}^2 \to \mathbb{C}^2 \otimes
\mathbb{C}^2$ assigns a Boltzmann weight to each combination of input
and output spins.

\subsubsection*{The five-vertex model}

For the KTW puzzle model, the nonzero weights are as follows.  We
write the vertex as $(a, b) \to (c, d)$, where $(a, b)$ are the
input spins (left, bottom) and $(c, d)$ are the output spins (right,
top):

\begin{center}
\renewcommand{\arraystretch}{1.3}
\begin{tabular}{ccccl}
\toprule
$a$ & $b$ & $c$ & $d$ & weight \\
\midrule
$0$ & $0$ & $0$ & $0$ & $1$ \\
$1$ & $1$ & $1$ & $1$ & $1$ \\
$0$ & $1$ & $1$ & $0$ & $1$ \\
$1$ & $0$ & $0$ & $1$ & $1$ \\
$0$ & $1$ & $0$ & $1$ & $1$ \\
\midrule
$1$ & $0$ & $1$ & $0$ & $0$ \\
\bottomrule
\end{tabular}
\end{center}

Five of the six possible ``ice-type'' vertex configurations (those
preserving the number of $1$s) have weight~$1$; the sixth---the
configuration $(1,0) \to (1,0)$, in which a spin-$1$ on the left
and a spin-$0$ from below exit unchanged to the right and above---is
\emph{forbidden} (weight~$0$).  This is the \emph{five-vertex model},
also known as the free-fermionic model.

In matrix form, writing $R$ in the basis $\{|00\rangle,
|01\rangle, |10\rangle, |11\rangle\}$:
\begin{equation}
\label{eq:rmatrix-puzzle}
  R \;=\;
  \begin{pmatrix}
    1 & 0 & 0 & 0 \\
    0 & 1 & 1 & 0 \\
    0 & 0 & 0 & 0 \\
    0 & 0 & 0 & 1
  \end{pmatrix}
  \quad\longleftrightarrow\quad
  \begin{pmatrix}
    1 & 0 & 0 & 0 \\
    0 & 0 & 1 & 0 \\
    0 & 1 & 1 & 0 \\
    0 & 0 & 0 & 1
  \end{pmatrix},
\end{equation}
where the two forms correspond to two natural orderings of the basis
(the right-hand matrix uses the ordering
$|00\rangle, |10\rangle, |01\rangle, |11\rangle$ to make the
block structure visible).  The zero row reflects the forbidden vertex.

The physical interpretation is conservation of ``charge'' (the number
of $1$s): every vertex has the same total spin on its input edges as
on its output edges.  In the puzzle picture, this is the constraint
that $1$-labeled edges form continuous paths through the triangle.

\subsubsection*{The spectral parameter}

The R-matrix above is the \emph{specialised} form, with all nonzero
weights equal to~$1$.  The general five-vertex R-matrix depends on a
spectral parameter $u$:
\begin{equation}
\label{eq:rmatrix-spectral}
  R(u) \;=\;
  \begin{pmatrix}
    1 & 0 & 0 & 0 \\
    0 & u & 1 & 0 \\
    0 & 0 & 0 & 0 \\
    0 & 0 & 0 & 1
  \end{pmatrix},
\end{equation}
where the weights of the five allowed vertices are $1, 1, 1, u, 1$
in the order listed in the table above (the ``straight-through''
vertex $(0,1) \to (0,1)$ acquires a weight $u$, while the crossing
vertex $(0,1) \to (1,0)$ retains weight~$1$).  The forbidden vertex
$(1,0) \to (1,0)$ remains forbidden at all values of $u$.

One verifies directly that $R(u)$ satisfies the Yang-Baxter equation
\eqref{eq:ybe} with additive spectral parameters:
\[
  R_{12}(u-v)\, R_{13}(u)\, R_{23}(v)
  \;=\;
  R_{23}(v)\, R_{13}(u)\, R_{12}(u-v).
\]
The puzzle model corresponds to $u = 0$ (or $u = 1$, depending on
normalisation conventions).  The spectral parameter does not change
\emph{which} configurations are allowed---it changes how they are
\emph{weighted}.  But because $R(u)$ satisfies the YBE for all~$u$,
the transfer matrices $T(u)$ form a commuting family, and all members
of the family share the same eigenvectors.  Changing $u$ changes the
eigenvalues, not the eigenspaces.  This is the power of the spectral
parameter: a single diagonalisation governs an entire family of
models.

\begin{remark}
The spectral parameter $u$ is precisely the variable appearing in the
transfer operator $T_{\mathrm{free}}(u)$ of Section~\ref{sec:classical-hopf}.
Setting $u = u_i$ in $R(u)$ and forming the transfer matrix via
\eqref{eq:transfer-matrix} recovers $T_{\mathrm{free}}(u_i)$ as an
operator on binary strings.  The formula
$s_{\lambda/\mu}(u_1, \ldots, u_n) = \langle\mu|
T_{\mathrm{free}}(u_1) \cdots T_{\mathrm{free}}(u_n)|\lambda\rangle$
is then a partition function with \emph{inhomogeneous} spectral
parameters, one per row.
\end{remark}

\subsubsection*{Connection to quantum groups}

The R-matrix \eqref{eq:rmatrix-spectral} is not an isolated
construction.  It is the image of the \emph{universal R-matrix} of
the quantum group $U_q(\widehat{\mathfrak{sl}}_2)$ in the fundamental
(two-dimensional) representation, at a specific value of $q$
corresponding to the free-fermionic point.  The Yang-Baxter equation
is the defining relation of the quantum group: any solution of the
YBE furnishes a representation of $U_q(\widehat{\mathfrak{sl}}_2)$,
and conversely, every representation yields an R-matrix satisfying the
YBE.

This connection is important for two reasons.  First, it places the
puzzle model in the vast landscape of quantum group representations,
where powerful algebraic tools (Bethe ansatz, fusion, quantum inverse
scattering) are available.  Second, it explains why R-matrices for
\emph{different} symmetric function families (Schur, Grothendieck,
Hall-Littlewood) are all solutions of the same equation: they are
different representations of the same quantum group, or of closely
related quantum groups in the same family.  We turn to these
extensions next.


\subsection{Extensions: the lattice model zoo}
\label{sec:lattice-extensions}

The five-vertex R-matrix of Section~\ref{sec:lattice-rmatrix}
governs the Schur case: its partition function computes LR
coefficients, and its transfer matrices build Schur functions row by
row.  But the Schur functions are only the first member of a
hierarchy of symmetric function families, each with its own
R-matrix and its own lattice model.  The integrable perspective
extends to each level of the hierarchy, and the R-matrix changes
in systematic, algebraically controlled ways.

\subsubsection*{Grothendieck polynomials}

Wheeler and Zinn-Justin \cite{wheeler-zinn-justin-2016} showed that
the \emph{stable Grothendieck polynomials}---the K-theoretic
analogues of Schur functions, representing Schubert classes in the
K-theory of the Grassmannian---arise as partition functions of a
\emph{trigonometric} five-vertex model.  The R-matrix is again
five-vertex (one vertex type forbidden), but the weights depend on an
additional parameter $\beta$ that deforms the Schur case:
\[
  R^{\mathrm{Groth}}(u; \beta) \;=\;
  \begin{pmatrix}
    1 & 0 & 0 & 0 \\
    0 & u + \beta & 1 & 0 \\
    0 & 0 & 0 & 0 \\
    0 & 0 & 0 & 1+\beta u
  \end{pmatrix}.
\]
At $\beta = 0$ this reduces to the Schur R-matrix
\eqref{eq:rmatrix-spectral}.  The parameter $\beta$ measures the
``K-theoretic deformation'': it introduces corrections beyond
cohomology, corresponding to higher-order terms in the Grothendieck
group.  The key fact is that this deformed R-matrix still satisfies
the Yang-Baxter equation, so the transfer matrices still commute, and
the partition function still computes the structure constants---now
the K-theoretic LR coefficients of Buch
\cite{buch-2002}.

\subsubsection*{Hall-Littlewood and spin models}

A different direction of generalisation leads to
\emph{higher-spin} models.  The five-vertex model has spins in
$\{0, 1\}$; allowing spins in $\{0, 1, \ldots, s\}$ for $s > 1$
produces lattice models whose partition functions are
\emph{Hall-Littlewood polynomials} and their relatives.

Borodin and Wheeler \cite{borodin-wheeler-2018} and
Wheeler and Zinn-Justin \cite{wheeler-zinn-justin-2025}
developed a family of \emph{coloured} and
\emph{higher-rank} vertex models whose partition functions compute
spin Hall-Littlewood polynomials.  Gunna, Wheeler, and Zinn-Justin
\cite{gunna-wheeler-zj-2025} showed that the structure constants of
these spin polynomials---analogues of LR coefficients for
Hall-Littlewood functions---are \emph{positive polynomials in~$q$}.
This positivity was not obvious from any previously known formula; it
emerged from the lattice model as a natural consequence of the vertex
weights being nonneg\-ative.

These higher-spin models are ``fully solvable'' in the sense of
satisfying the complete Yang-Baxter equation, in contrast to the
``quasi-solvable'' models of the next example.

\subsubsection*{Demazure characters and the $q=0$ specialisation}

Yang \cite{yang-2025} and Brubaker, Bump, Buciumas, and Gustafsson
\cite{brubaker-bump-buciumas-gustafsson-2021} constructed vertex
models whose partition functions are \emph{key polynomials}
(Demazure characters)---the characters of Demazure modules for
$GL_n$.  These models use a coloured five-vertex model at the
specialisation $q = 0$ of the Hecke algebra parameter.

At $q = 0$, the Hecke relation $(T_i - q)(T_i + 1) = 0$ degenerates
to $T_i(T_i + 1) = 0$, and the Demazure operators $\pi_i = T_i + 1$
become idempotent: $\pi_i^2 = \pi_i$.  The R-matrix still satisfies
the YBE restricted to pairs of adjacent rows (the ``column-level''
integrability of Buciumas and Scrimshaw), but the full three-line YBE
fails.  What survives is a weaker algebraic structure---a
\emph{coboundary} category governed by the cactus group rather than
the braid group.  We will develop this distinction systematically in
Section~\ref{sec:hierarchy}.

The states of the Demazure vertex model biject with
\emph{Demazure crystal} elements, and the partition function
manifestly computes the key polynomial.  Since the key polynomial
expands positively in the monomial basis, and the Schur function
expands positively in key polynomials (via the Kostka numbers), the
lattice model provides a direct bridge between the crystal
combinatorics of Kashiwara \cite{kashiwara-1990} and the integrable
lattice models of Zinn-Justin---a bridge that passes through the
Hecke algebra at $q = 0$.

\subsubsection*{The emerging pattern}

We can now see a pattern that is too regular to be coincidental.  At
each level of the character hierarchy---Schur, Grothendieck, key,
Hall-Littlewood, Macdonald---there is a solvable lattice model whose
partition function computes the corresponding symmetric (or
nonsymmetric) function:

\begin{center}
\renewcommand{\arraystretch}{1.3}
\begin{tabular}{lll}
\toprule
\textbf{Function} & \textbf{R-matrix type} & \textbf{Reference} \\
\midrule
Schur & rational 5-vertex & \cite{zinn-justin-2008} \\
Grothendieck & trigonometric 5-vertex & \cite{wheeler-zinn-justin-2016} \\
Key / Demazure & coloured 5-vertex, $q=0$ & \cite{yang-2025, brubaker-bump-buciumas-gustafsson-2021} \\
Hall-Littlewood & higher-spin & \cite{wheeler-zinn-justin-2025} \\
Macdonald & elliptic / higher-rank & --- \\
\bottomrule
\end{tabular}
\end{center}

Integrability is not an accident of the Schur case.  It is a
structural feature of the entire hierarchy of symmetric function
theories.  Each level has its own R-matrix, its own Yang-Baxter
equation (or a controlled degeneration thereof), and its own family of
commuting transfer matrices.

This raises the question that drives the remainder of the paper.  If
every character has a solvable model, what is the
\emph{relationship} between these models?  When we deform the Schur
R-matrix to the Grothendieck R-matrix, or specialise it to the
Demazure case, what algebraic structure governs the deformation?  Are
the bijections between different combinatorial models (puzzles,
tableaux, pipe dreams) for the \emph{same} character also controlled
by the Yang-Baxter equation?

The answer to the last question is yes, and the organising framework
is the \emph{Yang-Baxter groupoid} of Bump and Naprienko---the
subject of Section~\ref{sec:groupoid}.  The answer to the
deformation question leads to the \emph{categorical phase diagram} of
Section~\ref{sec:hierarchy}, where the vertical axis---braided,
coboundary, symmetric---classifies the integrability level, and the
horizontal axis---rational, trigonometric, elliptic---classifies the
deformation type.  Together, these two axes organise the full zoo of
lattice models into a single coherent picture.
